\introduction

\subsection{Motivation}
Industrial alloys, in particular steels and aluminium alloys, are selected for
mechanical components on the basis of their elastic properties: the Young's
modulus $E$ controls stiffness while the Poisson's ratio $\nu$ controls the
transverse response under load. Today these properties are obtained almost
exclusively from experiment, which is both expensive and time-consuming. If a
sufficiently accurate and inexpensive computational alternative existed, it
would dramatically accelerate the development of new alloy compositions for
specific mechanical demands.

Molecular dynamics (MD) simulations are the natural candidate. The reliability
of an MD prediction is, however, set entirely by the interatomic potential it
utilizes. Among the available functional forms, the Modified Embedded Atom Method
\citep{baskes1992,lee2001} is the de-facto standard for metallic and covalent
systems: it is computationally cheap relative to machine-learned potentials,
analytic, and physically interpretable.

\subsection{The Cross-Interaction Problem}
The fundamental obstacle to applying MEAM to industrial alloys is that
published MEAM parameter sets are fragmented. A typical paper provides a
3-4 element set, for example the Al--Mg--Zn potential of
\citet{Dickel_2018} or the Al--Si--Mg--Cu--Fe set of \citet{jelinek2012}, and
is silent on cross-interactions with elements that lie outside its scope.
A practically relevant alloy such as Al~7075 contains Al, Zn, Mg, Cu, Cr, Mn,
Fe, Si and Ti, and no single published MEAM file spans this composition.
Naively concatenating parameters from different libraries leaves the binary
cross-terms undefined and produces unphysical or unstable simulations.

Two routes around this problem have been pursued in the literature. The
first is manual fitting of new MEAM parameters against experimental and
\textit{ab initio} reference databases \citep{jelinek2012}, which is slow and
does not generalize to new element sets. The second is to abandon MEAM in
favour of fully data-driven machine-learning interatomic potentials
\citep{behler2007,shapeev2016,mishin2021}; these are flexible but require
large training sets and lose the analytic readability that makes MEAM useful
in industry.

\subsection{Approach Adopted in This Project}
The strategy adopted here is hybrid. Rather than replacing MEAM, neural
networks will eventually be used as \emph{surrogates} that optimize the MEAM
parameters themselves, starting from values seeded by DFT calculations. This
preserves the computational efficiency and transparency of the MEAM formalism
while exploiting DFT accuracy to fill the missing cross-interaction terms.
The four stages of the pipeline that are in scope for this interim
report---configuration generation, MD evaluation as baseline, DFT reference
calculations, and MEAM initialization---are described in detail in
Section~\ref{sec:progress}; a divergent composition-only neural-network
surrogate, used as a sanity check on the LAMMPS dataset, is described in
Section~\ref{sec:comp_nn}. The neural-network parameter optimization that
closes the loop is deferred to the second half of the semester.

\subsection{Research Question}
\textit{Can DFT reference calculations combined with neural-network surrogate
optimization produce multi-element MEAM interatomic potentials that accurately
reproduce the elastic properties of arbitrary metallic alloys, even when the
constituent elements do not share a common published MEAM parameterization?}
